\subsection{向量丛截面的积分}

在本号及下一号中，用 X 表示一个 \(C^{r}\) 类的实流形 \((1 \leq r \leq \infty)\)，且它局部有限维。

令 Y 为一个 \(C^{r}\) 类流形。若 \(r < \infty\)，考虑映射 \(f \mapsto j^{r}(f)\)，它从 \(\mathcal{C}^{r}(\mathrm{X};\mathrm{Y})\) 到 \(\mathcal{C}(\mathrm{X};\mathrm{J}^{r}(\mathrm{X},\mathrm{Y}))\)（《微分与解析流形，结果》，12.3.7）。\(\mathcal{C}(\mathrm{X};\mathrm{J}^{r}(\mathrm{X},\mathrm{Y}))\) 上的紧收敛拓扑在该映射下的逆像称为紧 \(C^{r}\) 收敛拓扑在 \(\mathcal{C}^{r}(\mathrm{X};\mathrm{Y})\) 上的拓扑；它是关于 K 的一致 \(C^{r}\) 收敛拓扑的上确界（《微分与解析流形，结果》，12.3.10），其中 K 遍历 X 的紧子集。

当 \(r = \infty\) 时，将紧 \(C^{\infty}\) 收敛拓扑定义在 \(\mathcal{C}^{\infty}(X;Y)\) 上，使其成为紧 \(C^{k}\) 收敛拓扑的上确界；换言之，它是使典则嵌入 \(\mathcal{C}^{\infty}(X;Y) \to \mathcal{C}^{k}(X;Y)\) 对所有 \(0 \leq k < \infty\) 均连续的最粗拓扑。

令 \(\mathrm{E}\) 为以 \(\mathrm{X}\) 为底空间的 \(\mathrm{C}^r\) 类实向量丛，令 \(\mathcal{S}^r(\mathrm{X};\mathrm{E})\) 为 \(\mathrm{E}\) 的 \(\mathrm{C}^r\) 类截面的向量空间。在本号中，赋予 \(\mathcal{S}^r(\mathrm{X};\mathrm{E})\) 由紧 \(\mathrm{C}^r\) 收敛拓扑在 \(\mathcal{C}^r(\mathrm{X};\mathrm{E})\) 上诱导的拓扑，也称为紧 \(\mathrm{C}^r\) 收敛拓扑；它使 \(\mathcal{S}^r (\mathrm{X};\mathrm{E})\) 成为完备、分离的局部凸拓扑向量空间（参见《微分与解析流形，结果》，15.3.1 以及《谱理论》（编写中））。

现在令 H 为李群，\(m: H \times X \to X\) 为 \(C^{r}\) 类的左作用律；置 \(hx = m(h, x)\)，其中 \(h \in H\)、\(x \in X\)。令 E 为以 X 为底空间的 \(C^{r}\) 类向量 H-丛（第 III 章，第 1 节，第 8 号，定义 4）。对于 \(s \in \mathcal{S}^{r}(X; E)\) 和 \(h \in H\)，将 \(^{h}s\) 记为 E 的截面 \(x \mapsto h.s(h^{-1}x)\)；映射 \((h, s) \mapsto ^{h}s\) 是 H 在空间 \(\mathcal{S}^{r}(X; E)\) 上的作用律。

引理 4. 作用律 \(\mathrm{H} \times \mathcal{S}^r (\mathrm{X};\mathrm{E}) \to \mathcal{S}^r (\mathrm{X};\mathrm{E})\) 是连续的。

根据 \(\mathcal{S}^{r}(\mathrm{X};\mathrm{E})\) 的拓扑定义以及《一般拓扑学》第 X 章第 3 节第 4 号定理 3，只需证明：对于任意整数 \(k \leq r\)，映射 \(f : H \times X \times \mathcal{S}^{k}(\mathrm{X};\mathrm{E}) \to \mathrm{J}^{k}(\mathrm{X};\mathrm{E})\) 满足 \(f(h,x,s) = j_{x}^{k}(^{h}s)\) 且是连续的。对于 \(h \in H\)，将 \(\tau_{h}\)（分别为 \(\theta_{h}\)）记为 X（分别为 E）上的自同构 \(x \mapsto hx\)。定义映射

\[
f _ {1}: \mathrm{H} \times \mathrm{X} \to \mathrm{J} ^ {k} (\mathrm{X}, \mathrm{X})
\]

\[
f _ {2}: \mathrm{H} \times \mathrm{E} \rightarrow \mathrm{J} ^ {k} (\mathrm{E}, \mathrm{E})
\]

\[
g: \mathrm{H} \times \mathrm{X} \times \mathcal {S} ^ {k} (\mathrm{X}; \mathrm{E}) \to \mathrm{J} ^ {k} (\mathrm{X}, \mathrm{E})
\]

由 \(f_{1}(h,x) = j_{x}^{k}(\tau_{h}),f_{2}(h,v) = j_{v}^{k}(\theta_{h}),g(h,x,s) = j_{hx}^{k}(s)\) 定义。我们有

\[
f (h, x, s) = f _ {2} (h, s (h ^ {- 1} x)) \circ g (h ^ {- 1}, x, s) \circ f _ {1} (h ^ {- 1}, x),
\]

因此，根据《微分与解析流形，结果》12.3.6，只需证明 \(f_{1}, f_{2}\) 和 \(g\) 连续。

现在 \(g\) 是复合映射

\[
\begin{array}{c c c} \mathrm{H} \times \mathrm{X} \times \mathcal {S} ^ {k} (\mathrm{X}; \mathrm{E}) & \xrightarrow {(m , \mathrm{Id})} & \mathrm{X} \times \mathcal {S} ^ {k} (\mathrm{X}; \mathrm{E}) \\ & \xrightarrow {(\mathrm{Id} , j ^ {k})} & \mathrm{X} \times \mathcal {C} (\mathrm{X}; \mathrm{J} ^ {k} (\mathrm{X}, \mathrm{E})) \quad \xrightarrow {\varepsilon} \quad \mathrm{J} ^ {k} (\mathrm{X}, \mathrm{E}) \end{array}
\]

其中 \(\varepsilon(x, u) = u(x)\)；由于映射 \(\varepsilon\) 连续（《一般拓扑学》第 X 章第 3 节第 4 号定理 3 的推论 1），\(g\) 连续。

令 \((h_{0}, x_{0}) \in \mathrm{H} \times \mathrm{X}\)；我们证明 \(f_{1}\) 在 \((h_{0}, x_{0})\) 处连续。存在图表 \((\mathrm{U}, \psi, \mathrm{F})\) 和 \((\mathrm{V}, \chi, \mathrm{F}')\)（均为 \(\mathrm{X}\) 的图表），以及开子集 \(\Omega\)（属于 \(\mathrm{H}\)），使得 \(x_{0} \in \mathrm{U}, h_{0} \in \Omega\) 且 \(m(\Omega \times \mathrm{U}) \subset \mathrm{V}\)。利用这些图表中的 \(\mathrm{J}^{k}(\mathrm{X}, \mathrm{X})\) 表达式，问题归结为证明：对于 \(1 \leq l \leq k\)，在 \((h_{0}, x_{0})\) 处映射 \((h, x) \mapsto \Delta_{x}^{l}(\tau_{h})\) 从 \(\Omega \times \mathrm{U}\) 到 \(\mathrm{P}_{l}(\mathrm{F}; \mathrm{F}')\) 连续，其中 \(\Delta_{x}^{l}(\tau_{h})(v) = \frac{1}{l!}\mathrm{D}^{l}\tau_{h}(x).v\) 对 \(v \in \mathrm{F}\) 成立（《微分与解析流形，结果》，12.2）。但是，\(\mathrm{D}^{l}\tau_{h}(x)\) 正是 \(l\) 阶偏导数 \(m(h, x)\) 关于 \(x\) 的，按假设它是连续的；因此 \(f_{1}\) 连续。\(f_{2}\) 连续性的证明类似，故引理得证。

命题 5. 假设群 H 是紧的，并将 H 上总质量为 1 的 Haar 测度记为 dh。令 s 为 E 的 \(C^{r}\) 类截面。

将 \(s^{\sharp}\) 记为向量积分 \(\int_{\mathrm{H}}{}^{h}s\,dh\)。则 \(s^{\sharp}\) 是 \(\mathrm{E}\) 的 \(\mathrm{C}^r\) 类截面，在 \(\mathrm{H}\) 下不变；对于 \(x \in \mathrm{X}\)，有 \(s^{\sharp}(x) = \int_{\mathrm{H}}hs(h^{-1}x)\,dh \in \mathrm{E}_x\)。将 \(s \mapsto s^{\sharp}\) 记为 \(\mathcal{S}^r(\mathrm{X};\mathrm{E})\) 上的自同态，它是到 H-不变截面子空间上的投影。

考虑映射 \(h \mapsto {}^{h}s\) 到 \(\mathcal{S}^{r}(\mathrm{X};\mathrm{E})\)；由引理 4 它是连续的。由于空间 \(\mathcal{S}^{r}(\mathrm{X};\mathrm{E})\) 分离且完备，积分 \(s^{\sharp} = \int_{H}{}^{h}s\,dh\) 属于 \(\mathcal{S}^{r}(\mathrm{X};\mathrm{E})\)（积分，第 III 章，第 3 节，第 3 号，推论 2）。线性映射 \(s \mapsto s(x)\) 从 \(\mathcal{S}^{r}(\mathrm{X};\mathrm{E})\) 到 \(E_{x}\) 且连续，因此 \(s^{\sharp}(x) = \int_{H}{}^{h}s(x)\,dh\) 对所有 \(x \in X\) 成立。显然 \(s^{\sharp}\) 在 H 下不变；若 s 是 H-不变截面，则 \(s^{\sharp} = s\)，故最后一个断言成立。

推论 1. 令 \(\mathrm{F}\) 为 Banach 空间，\(\rho : \mathrm{H} \to \mathbf{GL}(\mathrm{F})\) 为解析线性表示，\(f \in \mathcal{C}^r(\mathrm{X};\mathrm{F})\)。对于 \(x \in \mathrm{X}\)，置

\[
f ^ {\sharp} (x) = \int_ {\mathrm{H}} \rho (h). f (h ^ {- 1} x) d h.
\]

则 \(f^{\sharp}\) 是 \(\mathbf{C}^r\) 类的从 \(\mathbf{X}\) 到 \(\mathbf{F}\) 的态射，并且与 \(\mathbf{H}\) 的作用相容；对于 \(x \in \mathbf{X}\)，有（其中 \(\tau_h\) 表示自同构 \(x \mapsto hx\) 的 \(\mathbf{X}\)）

\[
d _ {x} f ^ {\sharp} = \int_ {\mathrm{H}} (\rho (h) \circ d _ {h ^ {- 1} x} f \circ \mathrm{T} _ {x} (\tau_ {h ^ {- 1}})) d h \in \mathscr {L} (\mathrm{T} _ {x} (\mathrm{X}); \mathrm{F}).\tag{14}
\]

第一断言由将命题应用于丛 \(X \times F\) 得到，其中赋予该丛以作用律 \((h; (x, f)) \mapsto (hx, \rho(h).f)\)。第二断言由积分，第 III 章第 3 节第 2 号命题 2 得到：对向量积分 \(f^{\sharp}\) 应用同态 \(d_{x} : \mathcal{C}^{r}(\mathrm{X};\mathrm{F}) \to \mathcal{L}(\mathrm{T}_{x}(\mathrm{X});\mathrm{F})\)，而该同态根据紧 \(C^{r}\) 收敛拓扑的定义是连续的。

推论 2. 令 F 为 Banach 空间，\(f \in \mathcal{C}^{r}(X;F)\)；置

\[
f ^ {\sharp} (x) = \int_ {\mathrm{H}} f (h x) d h
\]

对于 \(x \in \mathrm{X}\)。函数 \(f^{\sharp}\) 是 \(\mathbf{C}^r\) 类的，并且 \(f^{\sharp}(hx) = f^{\sharp}(x)\) 对于 \(x \in \mathrm{X}\)、\(h \in \mathrm{H}\) 成立。

推论 3. 令 F 为 Banach 空间，p 为满足 \(\geq 0\) 的整数，令 \(^{k}\Omega^{p}(\mathrm{X};\mathrm{F})\) 为 X 上取值于 F 且属于 \(C^{k}\) 类的 p 次微分形式空间（ \(2 \leq k + 1 \leq r\)）。对于 \(\omega \in {}^{k}\Omega^{p}(\mathrm{X};\mathrm{F})\)，置 \(\omega^{\sharp} = \int_{H} \tau_{h}^{*}\omega dh\)。则映射 \(\omega \mapsto \omega^{\sharp}\) 是 \(^{k}\Omega^{p}(\mathrm{X};\mathrm{F})\) 上的投影，其像为 H-不变形式子空间；有 \(d(\omega^{\sharp}) = (d\omega)^{\sharp}\) 对所有 \(\omega \in {}^{k}\Omega^{p}(\mathrm{X};\mathrm{F})\) 成立。

第一断言由将命题应用于向量 H-丛 \(\operatorname{Alt}^{p}(\mathrm{T}(\mathrm{X});\mathrm{F})\) 得到（第 III 章，第 1 节，第 8 号，例）。要证明第二断言，根据积分，第 III 章第 3 节第 2 号命题 2，只需证明映射 \(d:{}^{k}\varOmega^{p}(\mathrm{X};\mathrm{F})\to{}^{k-1}\varOmega^{p+1}(\mathrm{X};\mathrm{F})\) 是连续的，其中第一（分别为第二）空间赋予紧 \(C^{k}\) 收敛拓扑（分别为 \(C^{k-1}\) 收敛拓扑）。但这立即由利用半范数给出的这些拓扑的定义（《谱理论》（编写中））以及 d 是阶数 \(\leq 1\) 的微分算子这一事实得到（《微分与解析流形，结果》，14.4.2）。

